% !PRINTED 326
\weberpage{326}
Let us assume that in (16) $n$, $n'$ are themselves discriminants:
\[
n = D_1, \quad n' = D_2, \quad D_1 \equiv D_2 \pmod{D},
\]
then
\[
(D, D_1) = (D, D_2)
\]
and from (13) it follows that
\[
\begin{aligned}
(D, D_1) &= \pm(D_1, D) \\
(D, D_2) &= \pm(D_2, D),
\end{aligned}
\]
in which the upper signs hold when $D$ is positive, the lower ones when $D$, $D_1$, $D_2$ are all three negative, and different signs when $D$ is negative and $D_1$, $D_2$ are of different signs. Accordingly, by (16):
\weq{17}{(D_1, D) = \pm(D_2, D),}
in which the lower sign holds only when the \gesperrt{two numbers} $D$ and $D_1 D_2$ are both negative.

If $D_1$, $D_2$ are odd, we may set $D = 4m$, where $m$ is an arbitrary integer, and we obtain from (17)
\weq{18}{(D_1, m) = \pm(D_2, m),}
provided that
\weq{19}{D_1 \equiv D_2 \pmod{4m}}
\gesperrt{holds, the lower sign holding only when $m$ and $D_1 D_2$ are both negative.}

This theorem has so far been proved under the assumption that $D_1$ and $D_2$ are odd. But it holds in general. For the congruence (19) requires in the first place that $D_1$ and $D_2$ be at the same time even or at the same time odd. If they are both even, and $m$ is then even, (18) is correct, because both sides vanish. But if $m$ is odd and $D_1$, $D_2$ are even, then in (17) let us set $D = \pm m \equiv 1 \pmod{4}$, where $m$ is a positive odd number. Then from $D_1 \equiv D_2 \pmod{D}$ the congruence (19) follows, and by (17) we have
\[
(D_1, m) = (D_2, m),
\]
\weq{20}{\begin{aligned}
(D_1, m) &= (D_2, m),\\
(D_1, -m) &= \pm(D_2, -m), \\
(D_1, -1) &= \pm(D_2, -1),
\end{aligned}}
and consequently again
\[
(D_1, m) = (D_2, m).
\]
This formula therefore holds for every positive odd $m$, and by (20) then (18) is proved for a negative $m$ as well. The formulae (16) and (18) can be used for the calculation of the symbol $(D, n)$
%CONT>>
